Das $V(T)$-Diagramm zeigt vier Linien für eine feste Menge eines idealen Gases. Jede Linie ist eine Isobare ($p$ konstant), eine Isochore ($V$ konstant) oder eine Isotherme ($T$ konstant). Die Achsen sind in Einheiten von $V_0$ und $T_0 = \frac{p_0V_0}{nR}$ beschriftet.

\begin{center}
\begin{tikzpicture}[x=0.65cm, y=0.65cm, >=Stealth, font=\small]
 \draw[black!12] (0,0) grid[xstep=1, ystep=1] (8.001,8.001);
 \draw[->] (0,0) -- (8.6,0) node[right] {$T$ in $T_0$};
 \draw[->] (0,0) -- (0,8.6) node[above] {$V$ in $V_0$};
 \node[below left] at (0,0) {\small 0};
 \draw (1,0.08) -- (1,-0.08) node[below] {\small 1};
 \draw (2,0.08) -- (2,-0.08) node[below] {\small 2};
 \draw (3,0.08) -- (3,-0.08) node[below] {\small 3};
 \draw (4,0.08) -- (4,-0.08) node[below] {\small 4};
 \draw (5,0.08) -- (5,-0.08) node[below] {\small 5};
 \draw (6,0.08) -- (6,-0.08) node[below] {\small 6};
 \draw (7,0.08) -- (7,-0.08) node[below] {\small 7};
 \draw (8,0.08) -- (8,-0.08) node[below] {\small 8};
 \draw (0.08,1) -- (-0.08,1) node[left] {\small 1};
 \draw (0.08,2) -- (-0.08,2) node[left] {\small 2};
 \draw (0.08,3) -- (-0.08,3) node[left] {\small 3};
 \draw (0.08,4) -- (-0.08,4) node[left] {\small 4};
 \draw (0.08,5) -- (-0.08,5) node[left] {\small 5};
 \draw (0.08,6) -- (-0.08,6) node[left] {\small 6};
 \draw (0.08,7) -- (-0.08,7) node[left] {\small 7};
 \draw (0.08,8) -- (-0.08,8) node[left] {\small 8};
 \draw[very thick, blue] (0.3,5) -- (7.6,5);
 \node[draw, circle, inner sep=1pt, font=\footnotesize, right] at (7.6,5) {1};
 \draw[very thick, red] (0,0) -- (3.9,7.8);
 \node[draw, circle, inner sep=1pt, font=\footnotesize, above right] at (3.9,7.8) {2};
 \draw[very thick, teal] (6,0.3) -- (6,7.6);
 \node[draw, circle, inner sep=1pt, font=\footnotesize, above] at (6,7.6) {3};
 \draw[very thick, orange] (0,0) -- (8,4);
 \node[draw, circle, inner sep=1pt, font=\footnotesize, right] at (8,4) {4};
\end{tikzpicture}
\end{center}

\begin{abcliste}
\abc Welche Linie ist welche? Begründen Sie.
\abc Zwei Linien sind von derselben Art. Auf welcher der beiden ist der Druck grösser? Wie gross ist er auf jeder der beiden, in Einheiten von $p_0$?
\abc Welchen Wert hat die konstante Grösse auf den beiden anderen Linien?
\end{abcliste}
